\chapter*{Abstract}
\label{chap:Abstract}

La tesi analizza e confronta due tecnologie per il rilevamento della presenza umana in ambienti indoor, il sensore PIR (Passive InfraRed) e il radar mmWave a 24 GHz, per stabilire quale delle due rilevi una persona con maggiore affidabilità in uno scenario post-sisma. Il contesto applicativo in cui il confronto trova riscontro è il progetto DIPME, che prevede l’integrazione di sensori IoT (Internet of Things) negli arredi scolastici per segnalare, in seguito a un evento sismico, l’eventuale presenza di persone rifugiate sotto di essi.

Il nodo di quel progetto affida oggi il rilevamento della presenza a un sensore PIR. Tuttavia, a causa del suo principio di funzionamento, il PIR rileva il movimento ma non una persona immobile, che è proprio il caso di chi si è rifugiato sotto un arredo. Il lavoro valuta quindi l’impiego di un radar mmWave FMCW (Frequency Modulated Continuous Wave), per verificare se possa superare questo limite e quali siano i compromessi legati al suo utilizzo.

Nella prima parte vengono studiati il sensore PIR \pir{} e i due moduli radar \ldb{} e \ldc{}, con particolare attenzione ai loro principi di funzionamento, ai protocolli di comunicazione, ai parametri disponibili e ai limiti riportati nella documentazione dei costruttori. Per il confronto è stata realizzata una piattaforma di acquisizione basata su ESP32, in grado di campionare i sensori in modo sincrono e di accedere, nel caso dei radar, all’energia riflessa per ciascuna cella di distanza. A supporto della campagna è stata progettata e implementata un’interfaccia web, servita direttamente dall’ESP32, per la visualizzazione dei dati in tempo reale e la loro esportazione nello stesso formato.

La piattaforma è stata utilizzata per una campagna sperimentale che ha preso in esame la distanza di rilevamento, le latenze, i falsi positivi, l’attenuazione del segnale in presenza di ostacoli, la copertura angolare e il rilevamento del respiro. I risultati confermano che il PIR risponde al movimento e non alla persona ferma, mentre il radar l’ha rilevata in tutti gli scenari di confronto. Anche il secondo radar ha confermato questo comportamento, pur presentando alcuni limiti propri.

Sulle grandezze graduate fornite dal radar è stato inoltre sviluppato un indice di vitalità calcolato direttamente a bordo, che, oltre a indicare la presenza di una persona, fornisce una stima del suo livello di movimento. Il consumo energetico dei sensori è stato infine valutato sulla base dei dati forniti dai rispettivi costruttori.

%Il lavoro si conclude con la proposta di un’architettura ibrida in cui il PIR, grazie al suo basso consumo energetico, viene utilizzato per determinare quando attivare il radar